% -------------------------------------------------------------------------------------
\section{SCET\texorpdfstring{$_{\rm II}$}{II}、零箱减除与快度发散}\label{sec:scet-scet2}

有效理论的每个扇区都把圈积分扩展到全部动量，所以一个共线图在其动量趋于软（或超软）时会与软图重叠，重复计数必须减去：这就是零箱减除，它是 Collins 减除在有效理论中的对应。在 SCET$_{\rm I}$ 中共线与超软模式的虚度相差 $\lambda^2$；在 SCET$_{\rm II}$ 中共线与软模式的虚度相同，只在快度上分开，这时单个扇区的积分会出现维数正规化管不住的快度发散。本节依次讨论这两件事。

\begin{definition}[零箱减除]\label{def:scet-zero-bin}
(a) 设 $I_c(k)$ 是 SCET 共线图的被积函数（$k$ 为共线圈动量或实辐射动量）。其\emph{零箱}（zero-bin）$I_c^{(0)}(k)$ 是把 $k$ 按嵌套在其中的较小区域（软或超软模式）重新计数后 $I_c$ 的领头项；物理的共线贡献定义为 $\tilde I_c=\int(I_c-I_c^{(0)})$，多个嵌套区域时按包含关系逐级减除\cite{ManoharStewart2007}。这是\cref{def:method-subtraction}中 $C_R\Gamma=T_R(\Gamma-\sum_{R'<R}C_{R'}\Gamma)$ 在有效理论中的对应：$T_R$ 是共线近似，$R'$ 是嵌套于 $R$ 的软区域；它处理的正是共线圈积分与软圈积分的重叠。
(b) 红外发散也用维数正规化时，零箱减除与“除以软 Wilson 线矩阵元”的软减除等价，这在夸克形状因子与 $x_B\to1$ 的 DIS 中被显式比较，并以两圈阿贝尔的夸克形状因子作了检验\cite{IdilbiMehen2007}（两种形式的比较亦见\cite{Lee:2006nr}）；只有当零箱减除在每个共线因子中自洽地实施时，重复计数才被避免。
(c) 对推力的单圈喷注函数（\cref{sec:scet-rg}），实辐射零箱的被积函数 $\propto\delta^+(k^2)\,\frac{1}{n\cdot k\,\bar n\cdot k}\,\delta\big(n\cdot k-p^2/\bar n\cdot p\big)$（测量只约束 $n\cdot k$）。用 $\delta(k^2)=\delta(\bar n\cdot k\,n\cdot k-\bm k_{\mathrm T}^2)$ 积掉 $\bar n\cdot k$（雅可比因子 $1/n\cdot k$，且 $n\cdot k\,\bar n\cdot k=\bm k_{\mathrm T}^2$），剩下 $\frac{1}{n\cdot k}\int\dd^{2-2\varepsilon}k_\perp\,\bm k_{\mathrm T}^{-2}$：这是无标度积分，按\cref{lem:prelim-scaleless}为零。等价地，在 $\bar n\cdot k\to a\,\bar n\cdot k$、$k_\perp\to\sqrt a\,k_\perp$、$n\cdot k$ 固定之下，$\delta(k^2)$ 与 $1/\bar n\cdot k$ 各给出 $a^{-1}$，测度 $\dd(\bar n\cdot k)\,\dd^{2-2\varepsilon}k_\perp$ 给出 $a^{2-\varepsilon}$，零箱积分是 $a$ 的 $-\varepsilon$ 次齐次函数，只能为零。虚修正的零箱同样无标度。于是推力的减除后喷注函数等于未减除的喷注函数。% w18_s4 (D), w18_s6 (Z); App N prop:appN-thrust-zerobin is the Collins-form counterpart (pointer)
\uses{def:method-subtraction, lem:prelim-scaleless}
\end{definition}

在 Collins 减除的语言中，(c) 的对应陈述是\cref{prop:appN-thrust-zerobin}（[R, $R_{\mathrm{cond}}$]，只依赖于发散链项的取值约定）。

\paragraph{SCET$_{\rm II}$。}在横动量分布、$B$ 介子的遍举衰变等问题中，软模式 $(\lambda,\lambda,\lambda)$ 与共线模式相互作用，二者的虚度相同。

\begin{proposition}[SCET$_{\rm II}$ 与软 Wilson 线 \tagR]\label{prop:scet-scet2}
（$R_{\mathrm{PT}}$：在重整化微扰论中对红外正规化的部分子态严格。）（树图层次。）(a) 共线模式与软模式的虚度都是 $\lambda^2Q^2$；软动量 $k$ 被 $n$ 共线线吸收（或由其发射）后，中间线的虚度为 $(p+k)^2\simeq\bar n\cdot p\,n\cdot k\sim\lambda Q^2$，是硬共线的。(b) 在树图层次积掉这些硬共线线后，软胶子只通过算符中的软 Wilson 线
\begin{equation}\label{eq:scet-Sn}
S_n(x)=\mathrm P\exp\Big(\ii g\int_{-\infty}^0\dd s\;n\cdot A_s(x+sn)\Big)
\end{equation}
与 $n$ 共线夸克耦合（例如两喷注流化为 $\bar\chi_nS_n^\dagger\,\Gamma\,S_{\bar n}\chi_{\bar n}$），SCET$_{\rm II}$ 的领头阶拉氏量中没有软—共线耦合\cite{BauerEtAl2002,Hill:2002vw}。(c) 软贡献与共线贡献只在快度上分开，单个 SCET$_{\rm II}$ 积分可以有维数正规化不能正规化的\emph{快度发散}，必须引入快度正规子\cite[第~3、4~节]{ChiuEtAl2012b}。（本命题只涉及部分子振幅。）
\uses{def:scet-modes, prop:scet-wilson}
\end{proposition}

\begin{proof}
(a) 由 $2p\cdot k=\bar n\cdot p\,n\cdot k+n\cdot p\,\bar n\cdot k+2p_\perp\cdot k_\perp$，三项分别为 $\lambda$、$\lambda^3$、$\lambda^2$（以 $Q^2$ 为单位），而 $k^2\sim\lambda^2$。(b) 这与\cref{prop:scet-wilson}证明的后两步相同，只是 $n$ 与 $\bar n$ 互换：对出射的 $n$ 共线夸克，每一对“硬共线传播子 × 顶角”在领头幂次化为 $\bar u(p)\,\ii g\,\sla\epsilon^*t^a\,\frac{\ii\frac{\bar n\cdot p}{2}\sla n}{\bar n\cdot p\,(n\cdot k+\ii0)}=-g\,t^a\frac{n\cdot\epsilon^*}{n\cdot k+\ii0}\,\bar u(p)$（用 $\bar u_n\sla n=0$），$r$ 个胶子按次序求和给出 $S_n^\dagger$ 的展开（对 $n\cdot k\neq0$，$\pm\ii0$ 的取法不影响领头项；硬共线线也可以分两步积掉，先匹配到含硬共线与软模式的 SCET$_{\rm I}$，再匹配到 SCET$_{\rm II}$\cite{BauerEtAl2002,Hill:2002vw}）。对一个软胶子，精确发射振幅与程函因子 $-n\cdot\epsilon^*/n\cdot k$ 之差（对夸克自旋求和）的相对大小在一个代表性运动学点上为 $0.81$、$0.081$、$0.0081$（$\lambda=10^{-1},10^{-2},10^{-3}$），按 $\lambda$ 的一次方减小。% w18_s5 (H); c_scet quotes 0.78, 0.077, 0.0077 for a single spinor at the same point
软—共线耦合若出现在拉氏量中，就会把共线粒子推到硬共线虚度，故软胶子只能出现在算符的 $S_n$ 里。(c) 取一个软实辐射，其横动量 $\bm k_{\mathrm T}$ 被测量固定。令 $n\cdot k=k_T\ee^{-y}$、$\bar n\cdot k=k_T\ee^{y}$，
\begin{equation}\label{eq:scet-rapdiv}
\int\dd(n\cdot k)\,\dd(\bar n\cdot k)\,\frac{\delta(n\cdot k\,\bar n\cdot k-\bm k_{\mathrm T}^2)}{n\cdot k\,\bar n\cdot k}=\frac{1}{\bm k_{\mathrm T}^2}\int_{-\infty}^{\infty}\dd y ,
\end{equation}
在 $y\to\pm\infty$ 处发散；横向积分已被测量固定，$D$ 维测度对它不起作用。
\end{proof}

快度正规子有三类：Collins 的非类光 Wilson 线（\cref{def:tmd-collins}）；Chiu--Jain--Neill--Rothstein（CJNR）的 $\eta$ 正规子与快度重整化群\cite{ChiuEtAl2012a,ChiuEtAl2012b}，在软 Wilson 线的每个顶点插入 $w\,|2k^z/\nu|^{-\eta/2}$、在共线 Wilson 线的每个顶点插入 $w^2|\bar n\cdot k/\nu|^{-\eta}$（$w$ 为记账参数，最后取 $1$；先取 $\eta\to0$ 再取 $\varepsilon\to0$）\cite[式~(4.8)、(4.9)]{ChiuEtAl2012b}；Becher--Neubert 的解析正规子与“共线反常”\cite{BecherNeubert2011}。在 $\eta$ 正规子下，软圈的两个顶点合起来给出 $|2k^z|^{-\eta}=(2k_T|\sinh y|)^{-\eta}$，\cref{eq:scet-rapdiv}中的快度积分变为 $\int\dd y\,|2\sinh y|^{-\eta}=2^{-\eta}\Gamma(\tfrac\eta2)\Gamma(\tfrac12-\tfrac\eta2)/\sqrt\pi=\tfrac2\eta+\order{\eta}$：软辐射在正、负快度两侧都被截断。% w18_s5 (A2, A3)

\begin{proposition}[快度发散的处理与等价性 \tagA]\label{prop:scet-rapidity-rg}
对横动量依赖（TMD）的因子化：
(a) 快度重整化群（CJNR）、共线反常（Becher--Neubert）与 Collins 的非类光 Wilson 线（\cref{def:tmd-collins}）给出等价的 TMD 因子化\cite{Collins:2017oxh}：在每一种方案中，单个共线（束流）函数与软函数依赖于正规子，而进入截面的组合与之无关。
(b) 支配 $b$ 空间截面在固定 $b_T$ 下 $Q^2$ 依赖的演化核在三种方案中的对应为
\begin{equation}\label{eq:scet-kernels}
\tilde K(b_T;\mu)=\tfrac12\,\tilde\gamma_\nu^S(b_T;\mu)=-F_{q\bar q}(b_T^2,\mu)=-2D(b_T;\mu),
\end{equation}
其中 $\tilde K$ 是\cref{prop:tmd-evolution}的 Collins--Soper 核；$\tilde\gamma_\nu^S$ 由 $\nu\frac{\dd}{\dd\nu}\tilde S=\tilde\gamma^S_\nu\tilde S$ 定义（CJNR 的 $\eta$ 正规子，夸克 Wilson 线；束流函数的 $\tilde\gamma^B_\nu$ 满足 $\tilde\gamma^S_\nu+2\tilde\gamma^B_\nu=0$\cite[式~(5.38)]{ChiuEtAl2012b}）；$F_{q\bar q}$ 由 $[\tilde B_{q}\tilde B_{\bar q}]_{Q^2}=(b_T^2Q^2/b_0^2)^{-F_{q\bar q}}B_qB_{\bar q}$ 定义（Becher--Neubert\cite[式~(15)]{BecherNeubert2011}）；$D$ 是 $\delta$ 正规子（Echevarria--Idilbi--Scimemi）归一化中 $\partial\ln\tilde f/\partial\ln Q^2=-D$（$\zeta=Q^2$）的指数\cite[式~(5.4)]{EchevarriaIdilbiScimemi2012}。单圈时四者都等于 $-\frac{\as C_F}{\pi}L_b$，$L_b=\ln(\mu^2b_T^2/b_0^2)$。
（本命题陈述部分子层次的方案关系；强子层次的 TMD 因子化本身继承\cref{thm:tmd-factorization}的非微扰地位。）
\uses{prop:scet-scet2, def:tmd-collins, prop:tmd-evolution, thm:tmd-factorization}
\end{proposition}

\begin{proofsketch}
(1) 每种方案中束流函数与软函数的快度发散在乘积中相消（CJNR：$\tilde\gamma^S_\nu+2\tilde\gamma^B_\nu=0$\cite{ChiuEtAl2012b}；Becher--Neubert：$Q^2$ 依赖被完整地指数化\cite{BecherNeubert2011}；Collins：\cref{def:tmd-collins}中的平方根组合）。(2) 固定 $b_T$ 时 $b$ 空间截面 $H\tilde f_A\tilde f_B$ 对 $\ln Q^2$ 的导数是物理量：在 Collins 方案中取 $\zeta_A=\zeta_B=Q^2$，它等于 $\partial\ln H/\partial\ln Q^2+\tilde K$（\cref{prop:tmd-evolution}，[A]）；在 $\eta$ 方案中 $\tilde f=\tilde B\sqrt{\tilde S}$，$\tilde B$ 只通过 $\nu/\omega$ 依赖于大光锥动量 $\omega$，故 $\partial\ln\tilde f/\partial\ln\omega=-\tilde\gamma^B_\nu=\frac12\tilde\gamma^S_\nu$，而 $\zeta\propto\omega^2$；在另两种方案中乘积分别正比于 $(Q^2)^{-F_{q\bar q}}$ 与 $(Q^2)^{-2D}$。硬函数相同时比较即得\cref{eq:scet-kernels}；RG 一致性给出 $\frac{\dd\tilde K}{\dd\ln\mu}=-2\Gamma_{\rm cusp}$、$\frac{\dd F_{q\bar q}}{\dd\ln\mu}=2\Gamma_{\rm cusp}$\cite[式~(16)]{BecherNeubert2011}、$\frac{\dd D}{\dd\ln\mu}=\Gamma_{\rm cusp}$\cite[式~(5.14)]{EchevarriaIdilbiScimemi2012}。(3) 单圈核对：$\tilde K=-\frac{\as C_F}\pi L_b$（第~14~章）。在 $\eta$ 方案中，单个软实辐射的横动量被测量固定为 $\bm p$，由\cref{eq:scet-rapdiv}与上面的快度积分，$S^{(1)}(\bm p)=\frac{2g^2C_F}{(2\pi)^3\bm p^2}\big(\frac{\nu}{|\bm p|}\big)^{\eta}\big[\frac2\eta+\order{\eta}\big]$（$\varepsilon\to0$，$\int\dd^2p\,\delta^{(2)}(\bm p)=1$ 归一化），故 $\nu\frac{\dd}{\dd\nu}S^{(1)}=\frac{4\as C_F}{\pi}\mathcal L_0$，$\mathcal L_0=\frac1{2\pi}[1/\bm p^2]_+$ 在 $|\bm p|<\mu$ 内减除，与 CJNR 的单圈软函数（其式~(5.62)，换 $C_A\to C_F$）一致\cite{ChiuEtAl2012b}；由 $\int\dd^2p\,\ee^{\ii\bm p\cdot\bm b}\mathcal L_0=\int_0^\infty\frac{\dd p}p[J_0(pb)-\theta(\mu-p)]=-\frac12L_b$ 得 $\tilde\gamma^S_\nu=-\frac{2\as C_F}{\pi}L_b$（直接在 $b$ 空间计算 $\eta$ 正规化的单圈软函数，$\nu\frac{\dd}{\dd\nu}\tilde S^{(1)}=\frac{2\as C_F}{\pi}(-\frac1\varepsilon-L_b)$，去掉极点后结果相同）。Becher--Neubert 的单圈结果是 $F_{q\bar q}=\frac{C_F\as}\pi L_\perp$\cite[式~(40)]{BecherNeubert2011}，Echevarria--Idilbi--Scimemi 的是 $D=\frac{\as}{4\pi}\,2C_FL_T$\cite[式~(5.15)]{EchevarriaIdilbiScimemi2012}，其中 $L_\perp=L_T=L_b$。四者满足\cref{eq:scet-kernels}。% w18_s5 (B, C), w18_s6 (R1-R4)
完整论证见\cite{Collins:2017oxh,ChiuEtAl2012b,BecherNeubert2011,EchevarriaIdilbiScimemi2012}。
\end{proofsketch}

\paragraph{地位与缺口}
\cref{prop:scet-rapidity-rg}是全阶论证（[A]）。据该文摘要，文献\cite{Collins:2017oxh}由共线因子化与夸克形状因子的已知结果推出了 CSS 原始表述与 Collins 表述之间的系数关系，用于 $\order{\as^2}$、$\order{\as^3}$ 的已知计算，论证了从数据中提取的非微扰函数在两种表述中相同，并把它们与若干 SCET 表述联系起来：等价性是逐阶显式给出的，而不是对所有阶的证明。最弱的一步是 (a) 的“等价”本身：单个束流函数与软函数依赖于正规子，只有组合是物理的——快度发散标志着一个真实的附加因子化变量（$\zeta$ 或 $\nu$），而不是正规子的人为产物；这是方案之间的等价性问题，不涉及幂次误差。每一种方案都以 TMD 因子化（\cref{thm:tmd-factorization}，[A]）为前提，其瓶颈是\cref{asm:appN-tmd}、\cref{asm:appN-collins}与 Glauber 区的\cref{asm:appO-rewritten-cs,asm:appO-mixing}；TMD 算符的多重可重整性在文献中被论证而未被证明\cite[第~13~章]{Collins2011}。升级途径：在一个固定方案中把\cref{asm:appN-tmd}(T3) 证明到所有阶，并给出方案之间的有限重定义。
